\slajdid{10-s006}
\begin{frame}[label=10-s006]
\gusttekst\setlength{\abovedisplayskip}{3pt}\setlength{\belowdisplayskip}{3pt}\setlength{\abovedisplayshortskip}{2pt}\setlength{\belowdisplayshortskip}{2pt}\begin{columns}[c,onlytextwidth]\column{.28\textwidth}Zakočen\quad $V_{BE}<V_{\gamma T}$\column{.24\textwidth}\centering\begin{tikzpicture}[x=.85cm,y=.48cm,font=\sitneoznake,>=Latex]\ctikzset{bipoles/length=.65cm}\draw(-1,1)node[ocirc]{}node[left]{B}--(-.25,1);\draw(.35,1)--(1.1,1)node[ocirc]{}node[right]{C};\draw(0,-.3)--(0,-1)node[ocirc]{}node[below]{E};\draw[->](-.9,1.3)--(-.4,1.3)node[midway,above]{$I_B$};\draw[<-](.55,1.3)--(1,1.3)node[midway,above]{$I_C$};\draw[->](-.25,-.45)--(-.25,-.9)node[midway,left]{$I_E$};\draw(-.3,-.3)--(.3,-.3);\end{tikzpicture}
\column{.46\textwidth}
Struje $I_B=I_C=I_E=0$, dok su naponi $V_{CE}$ i $V_{BC}$ određeni spoljašnjim elementima sa kojima je tranzistor povezan.
\end{columns}\smallskip
\begin{columns}[c,onlytextwidth]\column{.28\textwidth}Aktivan režim\par\medskip $V_{BE}>V_{\gamma T}$ i $\beta_F I_B=I_C$\par\medskip $V_{BC}<V_{\gamma T}$\column{.24\textwidth}\centering\begin{tikzpicture}[x=.85cm,y=.48cm,font=\sitneoznake,>=Latex]\ctikzset{bipoles/length=.65cm}\draw(-1,1)node[ocirc]{}node[left]{B}--(-.25,1);\draw(.35,1)--(1.1,1)node[ocirc]{}node[right]{C};\draw(0,-.3)--(0,-1)node[ocirc]{}node[below]{E};\draw[->](-.9,1.3)--(-.4,1.3)node[midway,above]{$I_B$};\draw[<-](.55,1.3)--(1,1.3)node[midway,above]{$I_C$};\draw[->](-.25,-.45)--(-.25,-.9)node[midway,left]{$I_E$};\draw(-.25,1)--(-.25,.45);\draw(-.65,.45)--(.15,.45);\draw(-.5,.28)--(0,.28);\draw(-.25,.28)--(-.25,-.3)--(.4,-.3)--(.4,.2);\draw(.4,.4)circle(.2);\draw[->](.4,.55)--(.4,.25);\draw(.4,.6)--(.4,1);\node[left]at(-.65,.2){$V_{BE}$};\node at(-.5,.7){$+$};\node[right]at(.65,.4){$\beta_F I_B$};\end{tikzpicture}
\column{.46\textwidth}
$V_{BE}=0.6V$\par
Struja emitora je $I_E=(\beta_F+1)I_B$ dok su naponi $V_{CE}$ i $V_{BC}$ određeni spoljašnjim elementima sa kojima je tranzistor povezan
\end{columns}\smallskip
\begin{columns}[c,onlytextwidth]\column{.28\textwidth}Zasićenje\par\medskip $V_{BE}>V_{\gamma T}$ i $\beta_F I_B>I_C$\column{.24\textwidth}\centering\begin{tikzpicture}[x=.85cm,y=.48cm,font=\sitneoznake,>=Latex]\ctikzset{bipoles/length=.65cm}\draw(-1,1)node[ocirc]{}node[left]{B}--(-.25,1);\draw(.35,1)--(1.1,1)node[ocirc]{}node[right]{C};\draw(0,-.3)--(0,-1)node[ocirc]{}node[below]{E};\draw[->](-.9,1.3)--(-.4,1.3)node[midway,above]{$I_B$};\draw[<-](.55,1.3)--(1,1.3)node[midway,above]{$I_C$};\draw[->](-.25,-.45)--(-.25,-.9)node[midway,left]{$I_E$};\draw(-.25,1)--(-.45,1)--(-.45,.55);\draw(-.85,.55)--(-.05,.55);\draw(-.7,.38)--(-.2,.38);\draw(-.45,.38)--(-.45,-.3)--(.65,-.3)--(.65,.38);\draw(.25,.55)--(1.05,.55);\draw(.4,.38)--(.9,.38);\draw(.65,.55)--(.65,1);\node[left]at(-.85,.25){$V_{BES}$};\node[right]at(1.05,.25){$V_{CES}$};\node at(-.7,.8){$+$};\node at(.9,.8){$+$};\end{tikzpicture}
\column{.46\textwidth}
Struja kolektora i baze je određenja spoljašnjim elementima sa kojima je tranzistor povezan a struja emitora je kao i uvek $I_E=I_C+I_B$.
\end{columns}\medskip
Spoj baza emitor ponaša se kao idealan naponski izvor čija je vrednost napona u daljim primerima $V_{BES}=0.7V$. Napon između kolektora i emitora ćemo takođe smatrati konstantnim i smatrati da ne postoji unutrašnja otpornost između kolektora i emitora. Odnosno kao da se spoj baza kolektor ponaša kao idealan naponski izvor čija je vrednost napona u daljim primerima $V_{CES}=0.1V$
\end{frame}
